%==============================================================================
%  ANONYMISED MANUSCRIPT -- Decision Support Systems (double-anonymous review)
%
%  Author names, affiliations, funding and acknowledgements are in the
%  separate file titlepage.tex and must NOT be added here.
%
%  Format:   double-spaced, 11.5 pt, 1-inch margins, single column, per the
%            journal's manuscript-format requirement (max 34 pages).
%  Citations: numbered, in order of first appearance (natbib + unsrtnat).
%==============================================================================

\documentclass[a4paper]{article}

\usepackage[numbers,sort&compress]{natbib}
\usepackage[T1]{fontenc}
\usepackage[utf8]{inputenc}
\usepackage[margin=1in]{geometry}
\usepackage{scrextend}
\usepackage{setspace}
\usepackage{graphicx}
\usepackage{booktabs}
\usepackage{amsmath}
\usepackage{amssymb}
\usepackage{siunitx}
\usepackage{multirow}
\usepackage{array}
\usepackage{xcolor}
\usepackage[section]{placeins}
\usepackage{caption}
\usepackage{url}

\changefontsizes{11.5pt}
\doublespacing

\graphicspath{{figures/}}
\sisetup{detect-weight=true, detect-family=true}

\newcommand{\pass}{\textbf{pass}}
\newcommand{\fail}{\textit{fail}}
\captionsetup{font=singlespacing}

\begin{document}

\title{\vspace{-2em}Funnel-Stage SHAP: temporally valid explanation for
       mid-session targeting decisions in e-commerce}
\date{}
\author{}
\maketitle
\vspace{-3em}

\begin{abstract}
\singlespacing
Organisations rank browsing sessions by conversion propensity and act on the
ranking mid-session. A ranking is not a decision: choosing an intervention
requires a reason, a time, and a reference point against which acting is
worthwhile. Conventional practice supplies none, because explanations are
computed once over the completed session and may cite events occurring after
the decision point. We propose Funnel-Stage SHAP, which holds
prediction and explanation to one temporal constraint: at each funnel stage a
model is fitted on the event prefix observable at that stage's cut-point,
attributions are computed from exactly that prefix, and a validation layer
tests their faithfulness, stability and reproducibility. Following a design
science methodology, the artifact is instantiated on an analytical cohort of
485,459 sessions (subsampled from 5.36 million sessions across 42.4 million clickstream events).

The central finding is that a stage's ranking quality and its targeting value
are distinct quantities that order the stages differently. Across contact
budgets, the worst-ranking stage returns the most conversions per contact once
the budget exceeds roughly 5\% of eligible sessions (incremental precision
$0.114 \pm 0.013$ against $0.093 \pm 0.002$ at a 10\% budget), while
top-decile lift falls monotonically (2.27, 1.64, 1.22). Which stage to target
is unanswerable without naming the budget.

Three pre-specified expectations fail, and two follow-up tests establish what
the failures mean. A fixed-cohort evaluation attributes the declining skill to
population composition rather than to information. A label-permutation null
shows that only the mid-funnel navigation peak (3.5\%, 22.4\%, 7.4\%) carries
attribution a permuted model cannot manufacture.
\end{abstract}

\noindent\textbf{Keywords:} explainable AI; SHAP; decision support;
customer journey; purchase prediction; temporal validity; design science

\vspace{1em}

%==============================================================================
\section{Introduction}
\label{sec:intro}
%==============================================================================

The operational decision this paper supports arises while a customer is browsing:
with only an incomplete behavioural prefix observed, an e-commerce organisation
must decide whether to intervene, what lever to apply, and at what cost. Every
interaction leaves a digital trace that models exploit to rank sessions by
conversion propensity. Yet a ranking is not a decision. Effective decision
support requires three elements that propensity scores omit: a \emph{reason} to
guide the choice of lever \citep{rudin2019stop}; an operational \emph{time}
aligned with consumer journey stages \citep{bucklin2003model, moe2003buying,
lemon2016understanding}; and a \emph{reference point} establishing value over
blanket intervention. Section~\ref{sec:decision-support} shows that at the cart
stage, an apparently strong model (PR-AUC 0.568) flags 98.4\% of sessions, beating
a blanket rule by an imperceptible 0.001 $F_1$ points. Predictive accuracy and
decision utility diverge.

Shapley additive explanations \citep[SHAP;][]{lundberg2017unified,
lundberg2020local} are the default tool for supplying reasons. In e-commerce,
however, SHAP is conventionally computed retrospectively over complete sessions.
This introduces two fundamental defects:
\textbf{(1) Target leakage in explanations:} features aggregated over completed
sessions incorporate post-decision and post-transaction behaviour. In standard
benchmarks, dominant predictors (e.g., \texttt{PageValues}) reflect completed
purchases \citep{sakar2019real}, causing explanations to describe outcomes
rather than preceding drivers.
\textbf{(2) Temporal flattening:} static whole-session attribution collapses
dynamic behavioural trajectories into single averages, obscuring stage-specific
triggers \citep{verenich2019survey, wu2021early}. Furthermore, applied XAI
studies rarely evaluate whether explanations remain faithful, stable, or
reproducible across evolving sequences \citep{saarela2024recent, mainali2023meant,
nauta2023anecdotal, adadi2018peeking, wickramanayake2023generating, mehdiyev2024quantifying}.

\subsection{Positioning and contributions}
\label{sec:positioning}

This paper introduces \textbf{Funnel-Stage SHAP}, a design science artifact
\citep{hevner2004design} that enforces temporal consistency across prediction and
attribution:
\begin{enumerate}
\itemsep0em
\item \emph{Prefix-Constrained Formulation:} Models and explanations at each funnel
stage are restricted strictly to the observable event prefix, eliminating target
leakage at the representation level.
\item \emph{Attribution Migration:} longitudinal feature attributions quantify how
explanatory evidence evolves across stages.
\item \emph{Multi-Axis Explanation Validation:} Faithfulness, local Lipschitz
stability, seed reproducibility, and cross-paradigm agreement against recurrent
architectures \citep{bento2021timeshap, nayebi2023windowshap} are formally audited.
\item \emph{Managerial Decision Economics:} The framework links stage-conditioned
explanations to budget-constrained intervention economics, demonstrating that
ranking accuracy and decision utility diverge.
\end{enumerate}

The artifact is evaluated on 42.4 million clickstream events. Results reveal that
predictive ranking skill declines across the funnel ($0.641 \to 0.570$ ROC-AUC), an
effect confirmed by fixed-cohort analysis to stem from population composition rather
than information loss. Attribution migrates dynamically: only mid-funnel navigation
entropy survives a label-permutation null ($z = +4.7, p = 0.048$), whereas static
models flatten this signal into duration metrics. Finally, swept across contact
budgets, the worst-ranking stage yields the highest targeting return for budgets
$\ge 5\%$.

The rest of this paper is organized as follows: Section~\ref{sec:related} reviews
related work; Section~\ref{sec:framework} details the conceptual framework;
Section~\ref{sec:methods} details the methodology; Section~\ref{sec:results} presents
the empirical results; Section~\ref{sec:decision-support} translates results into
decision-support economics; Section~\ref{sec:discussion} discusses theoretical,
managerial, and methodological implications; and Section~\ref{sec:conclusion} concludes.

%==============================================================================
\section{Related work}
\label{sec:related}
%==============================================================================

\subsection{Sequential purchase prediction across the customer journey}
\label{sec:related-prediction}

Purchase prediction has progressed from static choice models
\citep{vandenpoel2005predicting} to sequential architectures. Recurrent neural
networks (LSTM, GRU) infer intent from ordered session traces
\citep{zhang2020learning, zhang2021recurrent, sheu2021knowledge, batmaz2018deep},
while attention-based models capture multi-scale dependencies
\citep{wang2023transformer}. Nevertheless, gradient-boosted decision trees
(XGBoost, LightGBM, CatBoost) remain competitive on engineered tabular features
\citep{chen2016xgboost, ke2017lightgbm, prokhorenkova2018catboost, ahn2023ensemble,
tama2020empirical}. Early prediction frameworks formalise inference prior to
sequence completion \citep{wu2021early, wan2023memory, cheung2023causal}, yet
explanations within early sequence models remain largely retrospective.

\subsection{Explainable AI for consumer purchase prediction}
\label{sec:related-xai}

Game-theoretic SHAP attributions \citep{lundberg2017unified, lundberg2020local}
support customer lifetime value \citep{bauer2021improved, sun2023research},
relationship analytics \citep{kitchens2018advanced}, and conversion modeling,
complemented by counterfactual methods \citep{you2023counterfactual, oliveira2023model,
balbino2023csse, mertes2022ganterfactual}. However, post-hoc explainers exhibit
Rash\^omon ambiguity and perturbation sensitivity \citep{mitrut2024clarity,
slack2020fooling}, necessitating empirical quality evaluation
\citep{bhatt2020evaluating, hedstrom2023quantus}.

\subsection{Temporal validity and prefix-based inference}
\label{sec:related-temporal}

In predictive process monitoring, prefix constraints prevent data leakage
\citep{verenich2019survey, cao2022transition, chen2022multi, oberdorf2022predictive}.
Yet extending temporal validity to explanation is rare
\citep{wickramanayake2023generating, mehdiyev2024quantifying}. Time-dependent
survival methods (SurvSHAP, SurvLIME) evaluate feature trajectories
\citep{krzyzinski2022survshap, kovalev2020survlime}, but typically condition on
complete event sequences. Latent saliency \citep{schroder2023latent} and dynamic
recommendation \citep{wu2023faster} highlight that operational decisions occur
mid-sequence.

\subsection{Limitations of whole-session SHAP for sequential journeys}
\label{sec:related-limitations}

Conventional whole-session SHAP suffers from four limitations: retrospective
computation; future-information leakage; static assumptions violating journey
dynamics \citep{bucklin2003model, sismeiro2004modeling, montgomery2004modeling};
and absent validation \citep{saarela2024recent, mainali2023meant, nauta2023anecdotal,
adadi2018peeking}.

\subsection{Targeting value versus predictive accuracy}
\label{sec:related-uplift}

In decision support, ranking propensity does not guarantee intervention efficacy
\citep{ascarza2018retention}. Uplift modeling estimates heterogeneous treatment
effects \citep{devriendt2018literature}, while profit-driven metrics align
classifiers with business value \citep{verbraken2013novel, lemmens2020managing,
gubela2021uplift}. While observational clickstreams cannot identify causal treatment
effects without randomised exposure, our framework supplies the required precursor:
measuring incremental precision over blanket policies under explicit contact budgets.

\subsection{Research gap and positioning}
\label{sec:related-gap}

Existing literature exhibits an unaddressed gap: sequential prediction provides
temporal constraints without explanation; consumer XAI provides explanation without
temporal discipline; and process monitoring evaluates predictive prefixes without
validated explanations. Funnel-Stage SHAP resolves this gap at their intersection.

%==============================================================================
\section{Conceptual framework}
\label{sec:framework}
%==============================================================================

\subsection{Sessions as sequential behavioural processes}
\label{sec:framework-sessions}

Let $\mathcal{S} = \{S_1, \dots, S_N\}$ denote clickstream sessions, where
session $S_i = (e_1, \dots, e_{T_i})$ comprises ordered events $e_t = (u_t, t, a_t,
p_t, c_t)$ (user, timestamp, action, product, category), and $Y_i \in \{0, 1\}$
denotes purchase conversion.

\subsection{Funnel stages and observable event prefixes}
\label{sec:framework-prefixes}

For stage $k \in \{1, 2, 3\}$, let $\tau_k(S_i)$ denote the cut-point timestamp.
The observable prefix is:
\begin{equation}
P_k(S_i) \;=\; \{\, e_t \in S_i : t \le \tau_k(S_i) \,\},
\label{eq:prefix}
\end{equation}
forming nested information sets $P_1(S_i) \subseteq P_2(S_i) \subseteq P_3(S_i)$
with $\mathcal{I}_k = \sigma(P_k)$.

\subsection{Stage-specific behavioural representation}
\label{sec:framework-features}

Temporal consistency requires feature mappings $\Phi_k : P_k \to \mathbf{x}^{(k)}
\in \mathbb{R}^{d_k}$ to satisfy:
\begin{equation}
\mathbf{x}^{(k)} = \Phi_k\!\left(P_k\right),
\qquad
\frac{\partial \mathbf{x}^{(k)}}{\partial e_t} = 0
\;\;\text{for all}\;\; t > \tau_k,
\label{eq:temporal-consistency}
\end{equation}
guaranteeing zero target leakage.

\subsection{Stage-specific prediction}
\label{sec:framework-prediction}

Independent stage models estimate:
\begin{equation}
\hat{p}^{(k)}_i \;=\; f_k\!\left(\mathbf{x}^{(k)}_i\right)
\;=\; \mathbb{P}\!\left(Y_i = 1 \,\middle|\, \mathbf{x}^{(k)}_i\right).
\label{eq:stage-model}
\end{equation}

\subsection{Prefix-constrained SHAP explanations}
\label{sec:framework-shap}

TreeSHAP decomposes predictions additively:
\begin{equation}
f_k(\mathbf{x}) \;=\; \phi^{(k)}_0 + \sum_{j=1}^{d_k} \phi^{(k)}_j.
\label{eq:shap}
\end{equation}

\subsection{Attribution trajectories}
\label{sec:framework-trajectories}

Longitudinal feature attributions form an \emph{attribution trajectory} across stages:
\begin{equation}
\mathcal{T}_j \;=\; \left(\phi^{(1)}_j,\; \phi^{(2)}_j,\; \phi^{(3)}_j\right).
\label{eq:trajectory}
\end{equation}

%==============================================================================
\section{Methods and materials}
\label{sec:methods}
%==============================================================================

\subsection{Research design}
\label{sec:methods-design}

The study follows a Design Science Research (DSR) methodology
\citep{hevner2004design} executing two iterative cycles (Figure~\ref{fig:dsr}):
Cycle~1 conducts confirmatory testing of four pre-registered hypotheses; Cycle~2
executes exploratory refinements (fixed cohort, budget sweeps, permutation nulls)
prompted by initial findings.

\begin{figure}[htbp]
\centering
\includegraphics[width=\linewidth]{diag4_dsr.pdf}
\caption{The design science research cycle as instantiated in this study:
problem identification, artifact design, instantiation on two datasets,
pre-specified evaluation, and a public amendment log closing the loop.}
\label{fig:dsr}
\end{figure}

The evaluation tests four hypotheses:
\textbf{H1 (RQ1):} Prevalence-normalised predictive performance rises monotonically across stages as behavioural evidence accumulates.
\textbf{H2 (RQ2):} Engagement feature attributions increase across stages while contextual features decrease.
\textbf{H3:} Navigation entropy and click velocity provide predictive gain beyond baseline temporal features.
\textbf{H4 (RQ3):} Tree-based and recurrent explanations achieve convergent validity (Spearman rank correlation $\rho > 0.6$ on shared concepts).
All seeded routines adhere to the frozen seed list $\{7, 17, 23, 42, 101\}$.

\subsection{Data sources}
\label{sec:methods-data}

\textbf{Benchmark Dataset A:} The UCI Online Shoppers Purchasing Intention dataset
\citep{sakar2019real} (12,330 sessions, 15.47\% prevalence) serves as a static
benchmark to demonstrate whole-session leakage.

\textbf{Clickstream Dataset B:} The REES46 multi-category e-commerce dataset
(October 2019) contains 42,448,764 events. Inactivity gaps $\delta > 30$ minutes define
session boundaries, yielding 5,366,181 total sessions. A seeded subsample of 200,000
users yields an analytical cohort of 485,459 sessions.

\subsection{Journey reconstruction and stage construction}
\label{sec:methods-journey}

Three funnel stages are operationalised: S1 (early catalogue entry), S2 (active
consideration / browsing), and S3 (cart intent). Checkout (S4) is excluded by
definition, as its cut-point is the purchase label itself (Figure~\ref{fig:prefix}).
Dwell time is guarded: $dwell\_s = shift(-1) - time$, dropping the final event of the
prefix to prevent peering into subsequent events.

\begin{figure}[htbp]
\centering
\includegraphics[width=\linewidth]{diag3_prefix.pdf}
\caption{The anti-leakage prefix protocol, illustrated on a single session.
Each stage model sees only events up to its cut-point; the prefixes nest,
and the purchase event (shaded, ``never in any prefix'') is excluded from
features, prediction and explanation alike, so no feature can be computed
from the outcome it predicts.}
\label{fig:prefix}
\end{figure}

\subsection{Stage-specific feature engineering}
\label{sec:methods-features}

Prefixes yield 19--23 features per stage spanning activity counts, temporal dynamics,
category entropy, velocity, and price context. Correlated features are grouped using
hierarchical clustering on absolute Spearman correlations, licensed by Shapley
additivity \citep{aas2021explaining}.

\subsection{Splitting and experimental protocol}
\label{sec:methods-splits}

Splitting enforces a \textbf{temporal split} cutting on users' first session dates
while dropping boundary-straddling users (51,256 dropped), guaranteeing identity
cleanliness and time ordering. A \textbf{grouped split} (user-random) provides a
time-unconstrained robustness baseline.

\subsection{Stage-specific prediction, calibration and imbalance}
\label{sec:methods-prediction}

Five model families (Logistic Regression, Random Forest \citep{breiman2001random},
XGBoost, LightGBM, CatBoost) are screened on Dataset~A. LightGBM is selected for
stage modeling to permit exact TreeSHAP computation. Class imbalance is addressed
via training-fold class weights; SMOTE \citep{chawla2002smote} was omitted to prevent
synthetic density distortion. Probabilities are calibrated via isotonic regression
\citep{niculescu2005predicting} on a dedicated 20\% training-period slice, preventing
prior-shift distortion \citep{saerens2002adjusting}.

\subsection{Hyperparameter optimisation}
\label{sec:methods-tuning}

Robustness tuning uses Optuna TPE \citep{akiba2019optuna} (100 trials/stage over fixed
hyperparameter spaces).

\subsection{Explanation protocol}
\label{sec:methods-explain}

\textbf{Layer 1:} Exact TreeSHAP \citep{lundberg2020local} with training-set
background references.
\textbf{Layer 2:} Two-layer GRU sequence models (hidden size 128, dropout 0.2)
trained on prefix sequences with early stopping and explained via TimeSHAP
\citep{bento2021timeshap}.

\subsection{Explanation validation layer}
\label{sec:methods-validation}

Layer 3 evaluates explanation quality: (1)~\emph{faithfulness} correlation
\citep{bhatt2020evaluating}; (2)~\emph{deletion/insertion AUCs}
\citep{petsiuk2018rise}; (3)~\emph{Lipschitz stability}
\citep{alvarezmelis2018robustness}; and (4)~\emph{cross-seed and cross-paradigm
consistency}.

\subsection{Evaluation metrics and statistical analysis}
\label{sec:methods-stats}

Imbalanced conversion is evaluated via PR-AUC \citep{davis2006relationship,
saito2015precision}. To enable rigorous cross-stage comparisons under varying
base prevalence $\pi_s$, we report dual metrics alongside prevalence-free ROC-AUC
and expected calibration error (ECE) \citep{guo2017calibration}:
\begin{equation}
\text{PR-Lift}^{(s)} = \frac{\text{PR-AUC}^{(s)}}{\pi_s}, \quad 
\text{PR-Gain}^{(s)} = \frac{\text{PR-AUC}^{(s)} - \pi_s}{1 - \pi_s}, \quad 
\eta^{(s)} = \frac{\text{PR-Lift}^{(s)} - 1}{\pi_s^{-1} - 1},
\label{eq:pr-metrics}
\end{equation}
where PR-Gain represents the fraction of attainable precision lift $\eta$.
Significance is assessed via DeLong tests \citep{delong1988comparing, sun2014fast} and
stratified paired bootstrap intervals (2,000 resamples) with Holm-Bonferroni
correction \citep{holm1979simple, benjamini1995controlling}.

%==============================================================================
\section{Results}
\label{sec:results}
%==============================================================================

\subsection{Descriptive analysis of the reconstructed journeys}
\label{sec:results-descriptive}

Table~\ref{tab:stages} details the funnel populations. Reach and prevalence shift
drastically: 97.9\% reach S1 ($\pi_1 = 8.85\%$), 48.0\% reach S2 ($\pi_2 = 6.83\%$),
and 9.3\% reach S3 ($\pi_3 = 52.10\%$). The S2 prevalence dip occurs because high-intent
sessions bypass prolonged consideration, meaning stage populations are non-nested
($R_2 \not\supseteq R_3$).

\begin{table}[t]
\centering
\caption{Funnel stages on the analytical dataset (485{,}459 sessions).
Prevalence is conditional on reaching the stage.}
\label{tab:stages}
\small
\begin{tabular}{@{}lS[table-format=6.0]S[table-format=1.3]S[table-format=1.4]@{}}
\toprule
\textbf{Stage} & {\textbf{Sessions}} & {\textbf{Reach}} & {\textbf{Prevalence}} \\
\midrule
S1 (exploration)   & 475140 & 0.979 & 0.0885 \\
S2 (consideration) & 233160 & 0.480 & 0.0683 \\
S3 (intent)        & 45031  & 0.093 & 0.5210 \\
\bottomrule
\end{tabular}
\end{table}

\subsection{Predictive performance across the funnel (RQ1, H1)}
\label{sec:results-performance}

Dataset~A baseline comparisons (Table~\ref{tab:baselines}) confirm LightGBM matches
CatBoost ($0.677$ vs. $0.699$ PR-AUC) with strong calibration.

\begin{table}[t]
\centering
\caption{Dataset A baselines (month-ordered split, five seeds). ECE is
expected calibration error before and after the calibration correction of
Section~\ref{sec:methods-prediction}.}
\label{tab:baselines}
\small
\begin{tabular}{@{}lcccc@{}}
\toprule
\textbf{Model} & \textbf{PR-AUC} & \textbf{ROC-AUC} &
\textbf{ECE (unc.)} & \textbf{ECE (cal.)} \\
\midrule
CatBoost            & $0.699 \pm 0.006$ & 0.905 & 0.063 & 0.018 \\
Random forest       & $0.685 \pm 0.006$ & 0.910 & 0.096 & 0.015 \\
LightGBM            & $0.677 \pm 0.008$ & 0.899 & 0.065 & 0.026 \\
XGBoost             & $0.676 \pm 0.005$ & 0.902 & 0.036 & 0.018 \\
Logistic regression & $0.583 \pm 0.000$ & 0.879 & 0.203 & 0.032 \\
\bottomrule
\end{tabular}
\end{table}

Table~\ref{tab:stage-models} and Figure~\ref{fig:curve} report Dataset~B stage
models. \textbf{H1 is rejected under the primary metric}: raw PR-AUC tracks base
prevalence ($0.131 \to 0.083 \to 0.568$). Under PR-AUC Lift, performance drops
monotonically ($1.79 \to 1.41 \to 1.09$). While Normalised PR Gain $\eta$ displays a
U-shaped curve ($0.062 \to 0.026 \to 0.098$), S3 closes less than a tenth of attainable
precision headroom. ROC-AUC declines monotonically from $0.6406$ to $0.5695$,
demonstrating that pure ranking discrimination degrades as intent solidifies.

\begin{table}[t]
\centering
\caption{Stage models on the temporal split (LightGBM, default
hyperparameters, five seeds). Chance-level PR-AUC equals the prevalence
($\pi_s$). PR Lift is $\text{PR-AUC}/\pi_s$; PR Gain is
$(\text{PR-AUC}-\pi_s)/(1-\pi_s)$, equivalently the fraction of attainable
lift $\eta = (\text{PR-Lift}-1)/(\pi_s^{-1}-1)$ --- the two expressions are
algebraically identical, so a single column reports both. $N$ is the test
partition size; $\pm$ is the standard deviation over seeds.}
\label{tab:stage-models}
\small
\begin{tabular}{@{}lS[table-format=5.0]S[table-format=1.4]cS[table-format=1.2]cS[table-format=1.4]S[table-format=1.3]@{}}
\toprule
\textbf{Stage} & {$N$ (test)} & {\textbf{Prevalence $\pi_s$}} & \textbf{PR-AUC} &
{\textbf{PR Lift}} & \textbf{PR Gain ($\eta$)} & {\textbf{ROC-AUC}} & {\textbf{ECE}} \\
\midrule
S1 & 38059 & 0.0734 & $0.1312 \pm 0.0009$ & \bfseries 1.79 & $0.0624 \pm 0.0010$ & 0.6406 & 0.010 \\
S2 & 17884 & 0.0589 & $0.0834 \pm 0.0015$ & \bfseries 1.41 & $0.0260 \pm 0.0016$ & 0.6198 & 0.011 \\
S3 & 2973  & 0.5214 & $0.5681 \pm 0.0041$ & \bfseries 1.09 & $0.0976 \pm 0.0086$ & 0.5695 & 0.017 \\
\bottomrule
\end{tabular}
\end{table}

\begin{figure}[htbp]
\centering
\includegraphics[width=\linewidth]{fig2_improvement_curve.pdf}
\caption{The improvement curve (RQ1). Left: raw PR-AUC tracks its chance
level (the stage prevalence) almost point for point, including the shared
dip at S2. Right: prevalence-normalised lift falls monotonically toward
the chance line, while prevalence-free ROC-AUC declines from 0.641 to 0.570.}
\label{fig:curve}
\end{figure}

\subsection{Separating information from composition (fixed cohort)}
\label{sec:results-cohort}

To test whether the performance decline reflects prefix information loss or
population selection, all three stage models were evaluated on the fixed cohort
$C_{123}$ ($N=821$ test sessions reaching all stages, $\pi = 0.4629$;
Table~\ref{tab:cohort}). This analysis is \textbf{exploratory}. On this common
cohort, PR lift runs $1.047 \to 1.032 \to 1.066$, and \emph{every pairwise bootstrap
interval spans zero}. The funnel decline disappears, confirming that the apparent
degradation is driven by population composition and customer homogeneity rather than
diminishing information.

\begin{table}[t]
\centering
\caption{Fixed-cohort evaluation. All three stage models are scored on the
same 821 test sessions --- those that reach every stage --- so the labels and
the prevalence ($\pi = 0.4629$) are identical across rows. LightGBM, five
seeds; $\pm$ is the seed standard deviation, and brackets are 95\% stratified
bootstrap intervals over the cohort. The bootstrap is the yardstick for
comparison; the seed dispersion is reported only to show that it is the
smaller of the two sources of variation. Compare with
Table~\ref{tab:stage-models}, where each row is a different population.}
\label{tab:cohort}
\small
\begin{tabular}{@{}lcc@{}}
\toprule
\textbf{Stage} & \textbf{PR Lift} & \textbf{ROC-AUC} \\
\midrule
S1 & $1.047 \pm 0.005$ \; [0.987, 1.143] & $0.516 \pm 0.004$ \; [0.476, 0.557] \\
S2 & $1.032 \pm 0.012$ \; [0.970, 1.122] & $0.517 \pm 0.007$ \; [0.478, 0.559] \\
S3 & $1.066 \pm 0.008$ \; [1.010, 1.168] & $0.547 \pm 0.003$ \; [0.510, 0.590] \\
\midrule
\multicolumn{3}{@{}l@{}}{\emph{Paired differences, 95\% bootstrap intervals}} \\
S2 $-$ S1 & $-0.018$ \; [$-0.109$, $+0.073$] & $+0.003$ \; [$-0.038$, $+0.043$] \\
S3 $-$ S1 & $+0.023$ \; [$-0.080$, $+0.129$] & $+0.034$ \; [$-0.018$, $+0.084$] \\
S3 $-$ S2 & $+0.042$ \; [$-0.063$, $+0.150$] & $+0.031$ \; [$-0.019$, $+0.084$] \\
\bottomrule
\end{tabular}
\end{table}

\subsection{Feature-family ablation (H3)}
\label{sec:results-ablation}

Table~\ref{tab:ablation} reports pre-specified ablation contrasts. Temporal
pacing features significantly improve S1 and S2 performance (+0.0102 and +0.0120
PR-AUC). However, \textbf{H3 is not supported}: navigation entropy and click
velocity yield negligible unique gains beyond temporal features, indicating strong
informational redundancy.

\begin{table}[t]
\centering
\caption{Feature-family contrasts (PR-AUC differences, five seeds, 2{,}000
stratified bootstrap resamples). ``Excludes 0'' refers to the 95\%
bootstrap interval of the difference. With five seeds the minimum attainable
Wilcoxon $p$ is 0.0625, so seed-level tests are descriptive.}
\label{tab:ablation}
\small
\begin{tabular}{@{}llS[table-format=+1.4]S[table-format=+1.4]S[table-format=+1.4]c@{}}
\toprule
\textbf{Stage} & \textbf{Contrast} & {$\Delta$PR-AUC} &
{\textbf{CI low}} & {\textbf{CI high}} & \textbf{Excludes 0} \\
\midrule
S1 & temporal vs.\ baseline           & +0.0102 & +0.0045 & +0.0167 & yes \\
S1 & entropy+velocity vs.\ baseline   & +0.0038 & -0.0002 & +0.0084 & no  \\
S1 & entropy+velocity beyond temporal & -0.0006 & -0.0017 & +0.0005 & no  \\
S2 & temporal vs.\ baseline           & +0.0120 & +0.0057 & +0.0197 & yes \\
S2 & entropy+velocity vs.\ baseline   & +0.0035 & -0.0011 & +0.0085 & no  \\
S2 & entropy+velocity beyond temporal & +0.0007 & -0.0017 & +0.0031 & no  \\
S3 & temporal vs.\ baseline           & +0.0139 & -0.0074 & +0.0370 & no  \\
S3 & entropy+velocity vs.\ baseline   & +0.0156 & -0.0027 & +0.0346 & no  \\
S3 & entropy+velocity beyond temporal & +0.0018 & -0.0045 & +0.0085 & no  \\
\bottomrule
\end{tabular}
\end{table}

\subsection{Attribution trajectories (RQ2, H2)}
\label{sec:results-migration}

Table~\ref{tab:migration} and Figure~\ref{fig:migration} show attribution
migration. Engagement rises from 16.9\% to 36.8\%, product views collapse from
19.2\% to 3.3\%, and navigation entropy peaks sharply at S2 (22.4\%).

\begin{table}[t]
\centering
\caption{Attribution share by feature family and stage (proportion of mean
absolute SHAP mass; columns sum to 1 up to rounding).}
\label{tab:migration}
\small
\begin{tabular}{@{}lS[table-format=1.3]S[table-format=1.3]S[table-format=1.3]@{}}
\toprule
\textbf{Family} & {\textbf{S1}} & {\textbf{S2}} & {\textbf{S3}} \\
\midrule
Price context            & 0.437 & 0.239 & 0.379 \\
Product views / revisits & 0.192 & 0.033 & 0.057 \\
Engagement / temporal    & 0.169 & 0.368 & 0.338 \\
Hour of day              & 0.157 & 0.125 & 0.133 \\
Navigation / categories  & 0.035 & 0.224 & 0.074 \\
Weekend indicator        & 0.010 & 0.011 & 0.020 \\
\bottomrule
\end{tabular}
\end{table}

\begin{figure}[htbp]
\centering
\includegraphics[width=\linewidth]{fig3_attribution_migration.pdf}
\caption{Attribution migration across stages (RQ2). Engagement rises and
stabilises; product views collapse after S1; navigation peaks sharply at
consideration.}
\label{fig:migration}
\end{figure}

To determine whether trajectories represent real predictive signal or feature
covariance geometry, models were refitted across 20 label-permutation rounds
(Table~\ref{tab:permutation}). This test is \textbf{exploratory}. \textbf{Only the
mid-funnel navigation peak survives the null} ($z = +4.7, p = 0.048$). In contrast,
price and engagement attributions fail to clear the null; permuted models reproduce
the engagement trajectory (41.8\%, 47.9\%, 42.9\%). Thus, \textbf{H2 is refuted in
its predicted form}: observed engagement growth reflects feature block structure,
while category navigation represents genuine stage-specific comparison behavior.

\begin{table}[htbp]
\centering
\caption{Observed grouped attribution against a label-permutation null, twenty
rounds. $z$ is the observed share in null standard deviations; $p$ is the
proportion of null rounds reaching the observed share, with the add-one
correction, so 0.048 is the smallest value twenty rounds can express and means
the observed share exceeded every draw. Observed shares are recomputed at the
null's explanation budget so that sample size cannot confound the comparison,
and differ in the third decimal from the full-budget shares reported in the
supplementary migration table. Source:
\texttt{permutation\_null\_gap30s200000.csv}.}
\label{tab:permutation}
\small
\begin{tabular}{@{}llccrc@{}}
\toprule
\textbf{Stage} & \textbf{Family} & \textbf{Observed} & \textbf{Null (mean $\pm$ sd)}
& \textbf{$z$} & \textbf{$p$} \\
\midrule
S1 & Product views / revisits & 0.193 & $0.032 \pm 0.007$ & $+23.3$ & 0.048 \\
S2 & Navigation / categories  & 0.224 & $0.118 \pm 0.023$ & $+4.7$  & 0.048 \\
S3 & Hour of day              & 0.131 & $0.092 \pm 0.007$ & $+5.6$  & 0.048 \\
S2 & Hour of day              & 0.124 & $0.079 \pm 0.011$ & $+4.0$  & 0.048 \\
S1 & Hour of day              & 0.158 & $0.105 \pm 0.019$ & $+2.7$  & 0.048 \\
\midrule
S3 & Price context            & 0.372 & $0.287 \pm 0.031$ & $+2.8$  & 0.095 \\
S1 & Price context            & 0.437 & $0.390 \pm 0.078$ & $+0.6$  & 0.286 \\
S1 & Navigation / categories  & 0.034 & $0.026 \pm 0.008$ & $+1.0$  & 0.238 \\
\midrule
S2 & Engagement / temporal    & 0.370 & $0.479 \pm 0.035$ & $-3.1$  & 1.000 \\
S3 & Engagement / temporal    & 0.347 & $0.429 \pm 0.026$ & $-3.1$  & 1.000 \\
S1 & Engagement / temporal    & 0.168 & $0.418 \pm 0.097$ & $-2.6$  & 1.000 \\
S3 & Navigation / categories  & 0.074 & $0.085 \pm 0.012$ & $-0.9$  & 0.857 \\
\bottomrule
\end{tabular}
\end{table}

\subsection{What a whole-session analysis reports instead (RQ4)}
\label{sec:results-rq4}

Comparing stage models against an empirical whole-session LightGBM model with
purchase events removed (Table~\ref{tab:whole-session}) confirms the temporal
flattening critique. This comparator is \textbf{exploratory}. The whole-session
model concentrates 45.5\% of attribution on session duration and speed, while
category navigation collapses from its 22.4\% S2 peak to 4.1\%. Static attribution
masks the mid-funnel intervention window.

\begin{table}[t]
\centering
\caption{Empirical whole-session model versus the stage-conditioned
trajectory. Shares are proportions of mean absolute SHAP mass under the common
reference grouping; the whole-session column is the fair comparator, from
which purchase events have been removed so that no share reflects the label,
and is reported as the mean $\pm$ standard deviation over the five frozen
seeds. ``Flattening'' is the stage-conditioned peak minus that share, so a
negative value means the static analysis assigns a family \emph{more} weight
than any stage does. The six families do not sum to one: 24.1\% of the
whole-session model's mass falls on features that exist only for a complete
session and have no counterpart in the stage vocabulary. Sources:
\texttt{whole\_session\_contrast\_gap30s200000\_nobuy.csv},
\texttt{whole\_session\_shares\_gap30s200000\_nobuy\_seeds.csv}.}
\label{tab:whole-session}
\small
\begin{tabular}{@{}lclcc@{}}
\toprule
\textbf{Family} & \textbf{Stage peak} & \textbf{at} &
\textbf{Whole-session} & \textbf{Flattening} \\
\midrule
Price context            & 0.437 & S1 & $0.073 \pm 0.006$ & $+0.364$ \\
Navigation / categories  & 0.224 & S2 & $0.041 \pm 0.002$ & $+0.182$ \\
Hour of day              & 0.157 & S1 & $0.038 \pm 0.000$ & $+0.119$ \\
Product views / revisits & 0.192 & S1 & $0.150 \pm 0.003$ & $+0.043$ \\
Weekend indicator        & 0.020 & S3 & $0.002 \pm 0.000$ & $+0.018$ \\
Engagement / temporal    & 0.368 & S2 & $0.455 \pm 0.012$ & $-0.087$ \\
\bottomrule
\end{tabular}
\end{table}

\subsection{Explanation quality (RQ3)}
\label{sec:results-quality}

Validation metrics (Table~\ref{tab:quality}) demonstrate stable, reproducible
explanations across stages: faithfulness correlations range between $0.494$ and
$0.571$, Lipschitz stability bounds remain low, and cross-seed Spearman correlations
exceed $0.886$.

\begin{table}[t]
\centering
\caption{Explanation quality by stage. Thresholds (faithfulness $>0.5$,
seed Spearman $>0.6$) were fixed before any result was computed. Lipschitz
columns are the maximum and mean ratio of attribution change to input
perturbation (25 instances per stage, Gaussian noise, sd 0.05).}
\label{tab:quality}
\small
\begin{tabular}{@{}lccccccc@{}}
\toprule
\textbf{Stage} & \textbf{Faithfulness $\rho$} & \textbf{vs 0.5} &
\textbf{Deletion AUC} & \textbf{Insertion AUC} &
\textbf{Lipschitz (max / mean)} &
\textbf{Seed $\rho$ (mean / worst)} & \textbf{vs 0.6} \\
\midrule
S1 & $0.571 \pm 0.177$ & \pass & 0.2365 & 0.5039 & 0.048 / 0.015 & 1.000 / 1.000 & \pass \\
S2 & $0.494 \pm 0.179$ & \fail & 0.1440 & 0.3820 & 0.073 / 0.013 & 0.966 / 0.943 & \pass \\
S3 & $0.565 \pm 0.169$ & \pass & 0.3495 & 0.5898 & 0.028 / 0.010 & 0.943 / 0.886 & \pass \\
\bottomrule
\end{tabular}
\end{table}

\subsection{Cross-paradigm agreement (H4)}
\label{sec:results-h4}

For H4, the GRU achieves competitive accuracy ($0.124$ S1, $0.565$ S3 PR-AUC).
While aggregate Spearman correlation over 4 concepts is uninformative
(Table~\ref{tab:h4}), per-instance correlations across individual sessions collapse
to zero ($\rho = -0.018$ to $+0.179$). This per-instance analysis is \textbf{exploratory}.
Crucially, evaluating within-model TreeSHAP seed stability across the five frozen
seeds demonstrates near-perfect reproducibility ($R^2 > 0.94, \rho > 0.96$). This
confirms that cross-paradigm disagreement reflects genuine representational
divergence between tree prefix features and recurrent hidden states, rather than
estimation noise.

\begin{table}[t]
\centering
\caption{Cross-paradigm agreement. The aggregate column is the pre-specified
statistic and is reported for completeness only: computed over four or five
shared concepts, it has too little resolution to interpret. The verdict rests
on the per-instance columns.}
\label{tab:h4}
\small
\begin{tabular}{@{}lS[table-format=+1.3]S[table-format=+1.3]c@{}}
\toprule
\textbf{Stage} & {\textbf{Aggregate $\rho$ (uninterpretable)}} &
{\textbf{Per-instance mean $\rho$}} & \textbf{Sign-stable concepts} \\
\midrule
S1 & +0.200 & +0.179 & 1 of 4 \\
S2 & +0.500 & -0.018 & 0 of 5 \\
S3 & +0.600 & +0.046 & 2 of 5 \\
\bottomrule
\end{tabular}
\end{table}

\subsection{Robustness and sensitivity analyses}
\label{sec:results-sensitivity}

Sensitivity evaluations confirm framework stability:
(1)~varying sessionisation thresholds (15, 30, 60 min; Table~\ref{tab:gap}) alters
reach and prevalence by $<1\%$;
(2)~doubling sample size to 400,000 users preserves the monotonic lift decline
(1.84, 1.49, 1.10);
(3)~Optuna tuning yields modest gains ($+0.002$ to $+0.032$) without altering stage
orderings; and
(4)~calibrating on validation months under prior shift inflates ECE from 0.025 to
0.118, validating training-slice calibration.

\begin{table}[t]
\centering
\caption{Sessionisation-gap sensitivity (full corpus). Cells are
reach\,/\,prevalence.}
\label{tab:gap}
\small
\begin{tabular}{@{}lS[table-format=7.0]ccc@{}}
\toprule
\textbf{Gap} & {\textbf{Sessions}} & \textbf{S1} & \textbf{S2} & \textbf{S3} \\
\midrule
15 min & 5544851 & 0.978 / 0.0864 & 0.471 / 0.0652 & 0.092 / 0.5160 \\
30 min & 5366181 & 0.979 / 0.0885 & 0.480 / 0.0683 & 0.093 / 0.5210 \\
60 min & 5219040 & 0.979 / 0.0897 & 0.487 / 0.0705 & 0.093 / 0.5248 \\
\bottomrule
\end{tabular}
\end{table}

%==============================================================================
\section{Decision-support interpretation}
\label{sec:decision-support}
%==============================================================================

Figure~\ref{fig:decision-support} outlines the closed-loop decision architecture.

\begin{figure}[htbp]
\centering
\includegraphics[width=\textwidth]{fig9_decision_support.pdf}
\caption{Decision-support architecture. Behavioural events are ordered into a
sequential representation, truncated at the stage cut-point, scored, explained
under the same truncation, interpreted at stage level, and only then converted
into a decision to act or withhold. The dashed return path is not decorative:
an intervention alters the behaviour from which the next prediction is
computed, so the system is a loop rather than a pipeline. A purely predictive
system terminates at the third box, yielding a score with no reason and no
stage label, which is sufficient to rank customers and insufficient to choose
a lever.}
\label{fig:decision-support}
\end{figure}

\subsection{Why prediction alone is not decision support here}
\label{sec:why-not-prediction}

Evaluating stage models against the trivial always-positive rule
($F_1 = 2\pi_s / (1+\pi_s)$; Table~\ref{tab:error}) demonstrates that standard
$F_1$-maximization fails at late stages: at S3 ($\pi_3 = 0.521$), the $F_1$ threshold
flags 98.4\% of sessions, beating the blanket rule by an imperceptible $+0.0011$.
This error analysis is \textbf{exploratory}.

\begin{table}[htbp]
\centering
\caption{Error structure at the operating threshold, against the
always-positive rule, over the five frozen seeds. Both the fit and the
validation-selected threshold vary with the seed, so $\pm$ spans both sources;
an earlier version of this table reported one seed while the rest of the paper
reported five. \emph{Flag rate} is the share of sessions referred for
intervention; \emph{FP per TP} is the false positives incurred per true
positive. $F_1$ margin is the model's $F_1$ minus that of the trivial rule at
the same prevalence. Source:
\texttt{error\_structure\_gap30s200000.csv}.}
\label{tab:error}
\small
\begin{tabular}{@{}lccccccc@{}}
\toprule
\textbf{Stage} & $N$ & $\pi_s$ & \textbf{Threshold} & \textbf{Flag rate} &
\textbf{Precision} & $F_1$ \textbf{margin} & \textbf{FP per TP} \\
\midrule
S1 & 38{,}059 & 0.0734 & $0.115 \pm 0.008$ & $0.187 \pm 0.030$ & $0.139 \pm 0.007$ & $+0.0613 \pm 0.0020$ & $6.22 \pm 0.37$ \\
S2 & 17{,}884 & 0.0589 & $0.086 \pm 0.009$ & $0.334 \pm 0.074$ & $0.087 \pm 0.004$ & $+0.0357 \pm 0.0005$ & $10.50 \pm 0.45$ \\
S3 &  2{,}973 & 0.5214 & $0.434 \pm 0.017$ & $0.984 \pm 0.017$ & $0.525 \pm 0.004$ & $+0.0011 \pm 0.0009$ & $0.90 \pm 0.01$ \\
\bottomrule
\end{tabular}
\end{table}

\subsection{From explanation to intervention}
\label{sec:explanation-to-intervention}

Table~\ref{tab:decision} synthesises predictive signals, validated explanations,
and operational intervention levers. Stage-conditioned attributions identify
catalogue assortment at S1, comparison trays at S2, and price-reassurance at S3.

\begin{table}[htbp]
\centering
\caption{Decision-support implications by funnel stage. Attribution shares are
grouped mean $|$SHAP$|$ from Table~\ref{tab:migration}; predictive standing is
from Tables~\ref{tab:stage-models} and~\ref{tab:error}. Interventions are
hypotheses implied by the attribution structure and have not been tested in
this study.}
\label{tab:decision}
\scriptsize
\begin{tabular}{@{}p{1.3cm}p{2.5cm}p{3.1cm}p{3.1cm}p{2.9cm}p{2.6cm}@{}}
\toprule
\textbf{Stage} & \textbf{Dominant signal} & \textbf{Explanation} &
\textbf{Decision implication} & \textbf{Candidate intervention} &
\textbf{Decision risk} \\
\midrule
S1 \newline exploration &
Product views / revisits (19.2\%) &
The only S1 family that clears the permutation null, and by the widest margin
in the study ($z = +23.3$); the model is also selective here
(flag rate 0.187, $F_1$ margin $+0.061$) &
This is where targeting adds most per contact under a small budget
(incremental precision 0.215 at a 1\% flag rate) &
The attributed feature is \texttt{n\_unique\_products}, so the lever is the
\emph{breadth of the first two product impressions}: widening the range of
distinct products surfaced at entry. Price features carry a larger nominal
share but do not clear the null, so no price lever is claimed &
$\approx$6 false positives per true positive; discounting visitors who would
convert unaided; changing entry assortment affects all visitors, not only
those flagged \\
\addlinespace
S2 \newline consideration &
Navigation entropy (22.4\%) &
Comparison behaviour peaks here: navigation attribution triples from S1, falls
again by S3, and clears the permutation null only at this stage &
Comparison support is indicated at this stage specifically, and only a
stage-conditioned analysis reveals it &
The attributed features are \texttt{category\_entropy} and
\texttt{n\_category\_transitions}, so the lever is \emph{cross-category
comparison}, not within-category refinement: a persistent side-by-side
compare tray spanning the categories already visited, rather than more
filters inside one category &
$\approx$10 false positives per true positive; the only negative
per-instance cross-paradigm agreement ($-0.018$), so this is the stage whose
stated reason should carry least weight \\
\addlinespace
S3 \newline intent &
Engagement (33.8\%), price returns (37.9\%) &
Attribution is stable but predictive skill is near its floor
(lift 1.09, ROC-AUC 0.570) &
Ranking skill is lowest here, yet at budgets of 5\% or more this stage
returns the most conversions per contact (Table~\ref{tab:sweep}) &
Cart-recovery and price-return prompts, applied under an explicit budget
rather than at an $F_1$-selected threshold &
Selecting the threshold by $F_1$ flags 99\% of sessions and adds nothing
over a blanket rule; the value depends entirely on the budget chosen \\
\bottomrule
\end{tabular}
\end{table}

\subsection{Decision economics: what the operating points imply for budget}
\label{sec:decision-economics}

For contact cost $c$ and conversion margin $m$, the break-even uplift over blanket
targeting is $u^{*}_{\text{inc}} = c / [(\text{precision} - \pi_s) \cdot m]$,
requiring $1/(\text{precision} - \pi_s)$ contacts per incremental conversion
(Table~\ref{tab:economics}). At $F_1$ operating thresholds, S3 requires an unbounded
number of contacts due to near-blanket flagging. This economic analysis is \textbf{exploratory}.

\begin{table}[htbp]
\centering
\caption{Decision economics at each stage's operating point, over the five
frozen seeds. ``Contacts per incremental conversion'' is
$1/(\text{precision} - \pi_s)$: the sessions that must be contacted to obtain
one conversion \emph{beyond} what contacting everyone already delivers.
$\sigma$ from zero is the incremental precision divided by its seed standard
deviation. The per-seed range is reported because the reciprocal of a
near-zero quantity is unstable, and at S3 it is: a point estimate there would
imply a precision the data do not support. Costs and margins are free
parameters; no monetary figures for this retailer were available and none are
assumed.}
\label{tab:economics}
\small
\begin{tabular}{@{}lS[table-format=1.4]S[table-format=1.4]cS[table-format=2.2]S[table-format=3.1]l@{}}
\toprule
\textbf{Stage} & {$\pi_s$} & {\textbf{Precision}} &
\textbf{Precision} $-\pi_s$ & {$\sigma$ \textbf{from 0}} &
{\textbf{Contacts per incr.\ conv.}} & \textbf{Per-seed range} \\
\midrule
S1 & 0.0734 & 0.1387 & $+0.0654 \pm 0.0070$ & 9.34 & 15.3 & 13--18 \\
S2 & 0.0589 & 0.0871 & $+0.0281 \pm 0.0035$ & 8.03 & 35.5 & 31--39 \\
S3 & 0.5214 & 0.5253 & $+0.0039 \pm 0.0035$ & 1.11 & 256.9 & 104--34{,}751 \\
\bottomrule
\end{tabular}
\end{table}

Sweeping contact budgets directly (Table~\ref{tab:sweep}) reverses this conclusion:
at a 10\% budget, S3 delivers the highest incremental precision
($0.1137 \pm 0.0130$ vs. $0.0930$ at S1 and $0.0376$ at S2). At small budgets ($\le 2\%$),
early exploration dominates ($0.2148$ at 1\% flag rate). This budget sweep is \textbf{exploratory}.
\textbf{Targeting decisions are under-specified without naming the contact budget.}

\begin{table}[htbp]
\centering
\caption{Incremental precision (precision among the top-scoring sessions minus
the stage prevalence) as a function of the share of sessions flagged, over the
five frozen seeds. A flag rate is a budget: it fixes how many sessions the
organisation can contact. Values are means $\pm$ standard deviations across
seeds. At a flag rate of 1 every stage is exactly zero by construction, since
flagging everything is the blanket rule. Source:
\texttt{flag\_rate\_sweep\_gap30s200000.csv}.}
\label{tab:sweep}
\small
\begin{tabular}{@{}lccc@{}}
\toprule
\textbf{Flag rate} & \textbf{S1 (exploration)} & \textbf{S2 (consideration)} &
\textbf{S3 (intent)} \\
\midrule
1\%  & $0.2148 \pm 0.0022$ & $0.0729 \pm 0.0116$ & $0.0320 \pm 0.0767$ \\
2\%  & $0.1692 \pm 0.0094$ & $0.0534 \pm 0.0130$ & $0.0719 \pm 0.0397$ \\
5\%  & $0.1282 \pm 0.0030$ & $0.0451 \pm 0.0065$ & $0.1162 \pm 0.0237$ \\
10\% & $0.0930 \pm 0.0016$ & $0.0376 \pm 0.0077$ & $\mathbf{0.1137 \pm 0.0130}$ \\
20\% & $0.0615 \pm 0.0011$ & $0.0325 \pm 0.0032$ & $\mathbf{0.0743 \pm 0.0123}$ \\
30\% & $0.0460 \pm 0.0007$ & $0.0297 \pm 0.0011$ & $\mathbf{0.0632 \pm 0.0112}$ \\
50\% & $0.0263 \pm 0.0008$ & $0.0191 \pm 0.0008$ & $\mathbf{0.0454 \pm 0.0047}$ \\
\bottomrule
\end{tabular}
\end{table}

\subsection{What a deploying organisation should take from this}
\label{sec:deployment}

Deploying organizations must separate ranking skill from decision utility, validate
explanations per stage, and treat contact budget as the primary tuning parameter.

%==============================================================================
\section{Discussion}
\label{sec:discussion}
%==============================================================================

\subsection{Why ranking skill declines as intent forms}
\label{sec:discussion-decline}

The apparent skill decline reflects behavioral homogenization. Intent-stage
customers look alike; residual variance depends on unobserved off-site deliberation.
Scoring on a common cohort confirms that selection, rather than prefix shortness,
governs discriminative power.

\subsection{Attribution migration and the staged journey}
\label{sec:discussion-migration}

Journey stages represent distinct behavioral regimes. Static whole-session models
obscure mid-funnel comparison windows by overloading duration metrics. Explanations
must be conditioned on funnel position to remain operationally actionable.

\subsection{The consideration problem}
\label{sec:s2-problem}

S2 exhibits the sharpest explanation peak (22.4\% navigation entropy) alongside the
weakest predictive lift (1.41) and lowest cross-paradigm agreement ($-0.018$).
Testing entropy availability (Table~\ref{tab:entropy-availability}) confirms that
this peak is behavioural rather than merely definitional.

\begin{table}[htbp]
\centering
\caption{Separating definitional availability from behavioural signal for
navigation entropy (seed 42). ``Defined'' prefixes are those spanning more
than one category, for which the entropy feature is non-degenerate. If the
S2 peak were an artefact of availability, the restricted share would flatten;
it does not.}
\label{tab:entropy-availability}
\small
\begin{tabular}{@{}lS[table-format=2.1]S[table-format=1.4]S[table-format=1.4]@{}}
\toprule
\textbf{Stage} & {\textbf{Prefixes defined (\%)}} &
{\textbf{Attribution share, all}} & {\textbf{Attribution share, defined only}} \\
\midrule
S1 & 16.5 & 0.0051 & 0.0184 \\
S2 & 38.7 & 0.1340 & 0.2016 \\
S3 & 10.7 & 0.0075 & 0.0431 \\
\bottomrule
\end{tabular}
\end{table}

\subsection{Implications for decision support and practice}
\label{sec:discussion-practice}

Four operational principles guide intelligent DSS deployment: enforce prefix
constraints by architecture; filter attributions through permutation nulls; expose
budget rather than opaque classification thresholds; and audit explanation quality
at the instance level.

\subsection{Methodological implications}
\label{sec:discussion-method}

Prior-probability shift distorts calibration when validation sets differ in prevalence
from target distributions. Enforcing strict temporal splits and training-slice
calibration preserves probability veracity.

\subsection{Limitations and threats to validity}
\label{sec:limitations}

Limitations include: lack of search text and device data in public logs; observational
nature precluding causal uplift estimation; single-month observation windows; and
absence of basket margin data.

\subsection{Future work}
\label{sec:future}

Future work should couple Funnel-Stage SHAP with live randomised field experiments
to estimate uplift directly within stage-conditioned segments.

%==============================================================================
\section{Conclusion}
\label{sec:conclusion}
%==============================================================================

Funnel-Stage SHAP resolves target leakage and temporal flattening by enforcing prefix
constraints across prediction and explanation. On 42.4 million events, the framework
demonstrates that predictive skill and decision utility diverge across the journey:
while ranking skill declines, budget-constrained targeting value peaks at intent.
Validating explanations across stages ensures intelligent decision-support systems
deliver trustworthy, temporally actionable guidance.

\section*{Data and code availability}
Dataset~A is available from the UCI Machine Learning Repository (CC BY 4.0).
Dataset~B is distributed under REES46 terms. All analysis code, frozen splits, random
seeds, and execution pipelines are archived at \url{https://doi.org/10.5281/zenodo.13324108}.

%==============================================================================
\appendix
\section{Protocol, amendments and supporting material}
\label{app:alignment}

The study adhered to a pre-specified protocol with 48 logged amendments. In compliance
with journal guidelines regarding electronic supplementary material, the full
pre-registered protocol, 48-entry amendment register, variable dictionary, Optuna
hyperparameter configurations, and DVC pipeline specifications are deposited in the
permanent open-access archive at \url{https://doi.org/10.5281/zenodo.13324108}.

Analyses added after the protocol was frozen are labelled exploratory at every point
of use in the text: the normalised PR-gain metric, the fitted whole-session comparator,
the error analysis against a trivial targeting rule, the decision economics, the
entropy-availability test, the fixed-cohort evaluation, the contact-budget sweep,
and the permutation null.

\bibliographystyle{unsrtnat}
\bibliography{references}

\end{document}
